% !TeX root = ../main.tex
\section{Kết Luận}
\label{sec:conclusion}

Bài báo đã trình bày Tactical Formation---một trò chơi cờ chiến thuật
trừu tượng mới kết hợp di chuyển kiểu Cờ Đam và bắt quân kiểu Reversi
trên bàn cờ hai lớp vị trí, hỗ trợ 2--4 người chơi. Luật chơi được
đặc tả đầy đủ và nhất quán cho cả ba biến thể. Mô hình toán học hình
thức cho chế độ 2 người làm rõ Định luật Bảo toàn 28 quân---phát biểu
hình thức cho khẳng định trực giác ``ăn = đổi phe, không loại bỏ''---phân
biệt Tactical Formation với cả Cờ Đam và Reversi.

AI đồng hành dựa trên tìm kiếm Minimax với cắt tỉa Alpha-Beta (biến thể
Paranoid cho chế độ nhiều người) chơi hợp lý trong ngân sách thời gian
tương tác và cho phép người chơi tự luyện tập khi không có đối thủ; giải
đấu bot-vs-bot xác nhận tìm kiếm sâu hơn thắng ổn định. Việc phân tích vì
sao MCTS với rollout ngẫu nhiên thuần không phù hợp cho lớp game
Reversi-like (rollout quá nhiễu) bổ sung cho công trình trước về
Othello~\citep{buro1999how}.

\paragraph{Hướng phát triển tiếp theo.}
\begin{enumerate}[label=(\roman*)]
  \item Mở rộng mô hình toán học hình thức sang chế độ 3 và 4 người,
        phân tích sự tồn tại của chiến lược cân bằng và động lực
        liên minh.
  \item Phân tích định lượng độ cân bằng game qua mô phỏng bot-vs-bot
        hàng trăm ván, ước lượng lợi thế tay đi đầu nếu có.
  \item Cải thiện hàm đánh giá AI với các đặc trưng chiến lược phức
        tạp hơn, hoặc học hàm đánh giá qua self-play theo hướng
        AlphaZero~\citep{silver2018alphazero}.
  \item Phát triển chế độ online để hỗ trợ nhiều người chơi từ xa.
  \item Khảo sát các biến thể luật chơi---ví dụ cho phép người đi chọn
        hướng sandwich thay vì bắt buộc lấy hướng nhiều nhất---và
        đánh giá ảnh hưởng đến chiều sâu chiến thuật.
  \item Mô tả Tactical Formation trong một hệ general game system như
        Ludii~\citep{piette2020ludii} để cộng đồng có thể chơi tự động
        và so sánh với các game cờ trừu tượng khác.
\end{enumerate}
